\section{Introduction}

Code obfuscation is widely used to make software difficult to inspect
while preserving its intended functionality. In benign settings,
developers employ obfuscation to protect intellectual property,
implementation details, and proprietary business logic. In adversarial
settings, however, malware authors routinely exploit the same
techniques to conceal malicious behavior and frustrate program
analysis and reverse engineering. Common transformations include
identifier renaming, control-flow restructuring, dead-code insertion,
expression rewriting, call indirection, etc. Although such
transformations are intended to preserve program semantics, they can
substantially obscure how those semantics are expressed in the source
code, thereby increasing the difficulty of understanding the
behavior of malware code~\cite{collberg1997taxonomy,collberg2002manufacturing,
Obfuscation-The-Hidden-Malware,
An-Observational-Investigation-of-Reverse-Engineers-Process-and-Mental-Models}.

This difficulty is well established for human analysts. Nguyen
{\em et al.}~\cite{nguyen2026effectcodeobfuscationhuman}, for example,
showed that code obfuscation significantly reduces humans' accuracy in
reasoning about program outputs. Such findings are particularly
important for malware analysis, where analysts must often infer the
behavior of code whose surface structure has been deliberately
engineered to resist comprehension.

Large language models (LLMs) offer a promising alternative for
automating such reasoning. Models have demonstrated strong
capabilities in code generation, summarization, vulnerability
detection, and program understanding
\cite{chen2021evaluatinglargelanguagemodels,li2022codebert,
guo2021graphcodebert,roziere2023codegen,
Competition-level-code-generation-with-AlphaCode}.
However, most evaluations are performed on naturally written programs,
where meaningful identifiers, familiar control structures, and common
coding idioms provide abundant lexical and structural cues. Obfuscated
programs remove or distort many of these cues while retaining the
underlying computation, creating a useful stress test of whether an
LLM actually reasons about program behavior or instead depends on
surface regularities.

Recent studies suggest that this distinction matters.
% >>> MISSING REFERENCE, 2026-09-14. `ase26` has no entry in references.bib or
% references-codesteer.bib, and no match anywhere in papers/ or the sibling projects --
% the only Rong in the bibliography is rong2026codesteer, which is attention steering, not
% obfuscation. The build reports `Citation \`ase26' undefined`. Not invented here: the
% sentence makes a specific empirical claim and needs the real paper. <<<
Rong {\em et al.}~\cite{ase26} showed that semantics-preserving
obfuscations can cause substantial degradation in LLM program
reasoning, including program-output prediction. Similarly, Nikiema
{\em et al.}~\cite{nikiema2025codebarrierllmsactually} reported a
statistically significant decline in model performance as the
complexity of obfuscation increases. Together, these findings indicate
that this limitation is especially problematic
for security applications: an AI-assisted malware analyzer should not
cease to understand a malicious computation merely due to obfuscation.

%today's code LLMs are not invariant to transformations that leave program semantics unchanged.

An important challenge, however, has received considerably less
attention: \emph{real obfuscated programs rarely contain only one
transformation}. Obfuscators can combine multiple techniques so that,
for example, renamed identifiers are embedded inside flattened control
flow, rewritten expressions, and additional irrelevant code. Such
\emph{stacked obfuscation} creates a fundamentally harder deployment
setting. Success on each transformation independently does not
necessarily imply success when those transformations occur together.
Consequently, simply demonstrating that a model can be adapted to
individual obfuscations is insufficient. For practical malware
analysis, the learned capability must also \emph{compose}.

This observation naturally suggests several ways to compose knowledge
learned from individual transformations. One could train a specialist
for each obfuscation and use a \emph{router} to select the appropriate
specialist at inference time. Alternatively, one could combine the
specialists in parameter space through \emph{model merging}. A third,
and seemingly more direct, strategy is to increase \emph{training
breadth} by fine-tuning a single model on examples from many individual
obfuscations, hoping that knowledge learned separately will recombine
when transformations are stacked. We evaluate these three strategies, together with a
clean-code control and a paired-consistency objective that anchors the
student to its own behaviour on the unobfuscated parent, across eight
models, and we evaluate them in three regimes rather than one:
compositions of the transformations seen in training, an obfuscator
family held out entirely, and the same programs queried in the reverse
direction. The regimes disagree, and the disagreement is the result.
On seen compositions every adaptation works and the choice between
them is worth at most five points. Outside them the ordering inverts.
Training breadth --- the obvious way to spend the data --- is the worst
performer on the unseen family, below every arm assembled from
per-transform specialists on all eight models and below an adapter
trained on unobfuscated code alone on six. Broader exposure improves
recombination within the training distribution while making the model
\emph{less} robust outside it.

Our analysis reveals why. The apparent robustness obtained from
obfuscation-specific adaptation does not primarily arise from learning
the semantics that remain invariant under transformation. Instead,
models learn regularities associated with the \emph{surface artifacts}
left by particular obfuscators. Evidence for this behavior appears
across multiple analyses: adaptation transfers poorly between distinct
mechanisms within the same obfuscation family; the benefit of broad
training is substantially stronger when recognizable identifier cues
remain available; and improvements disappear when the same programs
are queried in a different reasoning direction. Thus, stacking exposes
a deeper problem than merely insufficient training coverage:
obfuscation-specific fine-tuning can teach the model how particular
transformations \emph{look} without teaching it to consistently recover
the behavior that those transformations preserve.

This diagnosis suggests a different principle for robust adaptation.
Rather than attempting to compose transformation-specific knowledge,
we anchor learning to something that is invariant across all such
transformations: the model's behavior on the corresponding
\emph{unobfuscated program}. For every obfuscated training program
$x_{\mathrm{obf}}$, we retain its clean parent
$x_{\mathrm{clean}}$. A teacher model, tuned for reasoning over clean
code, processes $x_{\mathrm{clean}}$, while the student processes
$x_{\mathrm{obf}}$. In addition to the ordinary task loss, we minimize
the KL divergence between their output distributions:
%
\begin{equation}
\mathcal{L}
=
\mathrm{CE}\!\left(y \mid x_{\mathrm{obf}}\right)
+
\lambda\,
D_{\mathrm{KL}}\!\left(
p_{\mathrm{teacher}}(\cdot \mid x_{\mathrm{clean}})
\,\Vert\,
p_{\mathrm{student}}(\cdot \mid x_{\mathrm{obf}})
\right).
\label{eq:paired_consistency}
\end{equation}
%
The first term exposes the student to obfuscated code, whereas the
second explicitly encourages it to preserve the behavioral response of
a model reasoning over the semantics-equivalent clean program.
Importantly, this objective does not require recognizing which
obfuscations are present, selecting among specialists, merging
transformation-specific parameters, or recovering the original source
code. Instead, it
changes the learning target itself: rather than anchoring adaptation to
the surface patterns of individual transformations, it anchors the
student to the clean-code behavior that those transformations should
preserve.

Across the two out-of-distribution regimes a single arm behaves
differently from the rest. Weight-space merging of the per-transform
specialists gives up three to six points on seen compositions, where
fitting the training distribution pays, and gives up nothing anywhere
else: it is above the untuned model on the reverse task on seven of
eight models and is never significantly below it in any regime on any
model. The mechanism is visible in what the other arms do when they
fail backwards. Their failures are not malformed answers but the gold
\emph{forward} answer --- the program's return value, correctly
computed, offered in reply to a question that asked for an input. Every
single-adapter arm does this at a measurable rate; both merges sit at
the floor of it on all eight models. Monolithic training fits the task
as it was posed, these transformations and this direction, and
composing separately-trained specialists in weight space does not
produce that fit.

We state the bound plainly. Merging is an existing technique, applied
here without modification; the contribution is not the algorithm but
the evaluation that separates the arms and the finding that
composition survives what monolithic training does not. Merging does
not create a new ability to read obfuscation it has never seen --- on
the unseen family it is competitive with, not superior to, the mixture
of the same specialists --- and it is not the best forward arm
anywhere. What it does is avoid paying, outside the training
distribution, for a fit inside it.

{\em Contributions.} This paper makes the following main contributions:

\textbf{1. A three-regime evaluation of obfuscation adaptation, and the disagreement between the
regimes.} We evaluate six adaptations on eight models across four lineages in three regimes:
compositions of seen transformations, a held-out obfuscator family, and the reverse task (given a
return value, produce a call that yields it). On seen compositions the adaptations are within five
points of one another; on the other two the ordering inverts. Reporting only the first regime,
as prior work does, supports the opposite conclusion to reporting all three.

\textbf{2. Monolithic training is the worst use of the data outside the training distribution.} A
single adapter trained on the pooled data of all six training conditions is below every arm assembled
from per-condition specialists on the unseen family, on all eight models --- and below an adapter trained on unobfuscated
code alone on six. The comparison is data-matched by construction: the composed arms are built from
specialists trained on exactly the rows breadth trains on, at the same rank, from the same base model.

\textbf{3. Adaptation locks the direction it was trained in, and we measure it directly.} Every
single-adapter arm loses backward accuracy on the same programs it gained forward accuracy on. The
failures are not parse errors: the most common ungradable backward reply is exactly the gold forward
answer. We report this collapse rate per arm per model; the untuned models and both merges sit at its
floor, and the highest rate belongs to the objective whose extra loss term most directly optimises for
the forward answer.

\textbf{4. Weight-space merging as the arm that never breaks, with its price stated.} Merging the same
specialists is above the untuned model on the reverse task on seven of eight models and is never
significantly below it in any regime on any model, at a cost of three to six forward points where
fitting the training distribution pays, one adapter, and no inference-time machinery. It is an existing
technique used unmodified; the claim is about what the evaluation reveals, not about the algorithm.
